
\begin{frame}[fragile]{Récapitulatif SQL: le bloc \texttt{select-from-where}}

Dans cette session: présentation de quelques extensions pratiques au SQL fondamental.
\vfill

\begin{redbullet}
  \item Les valeurs nulles
  \item La jointure externe
  \item Le tri
\end{redbullet}

\vfill

\alert{Avant d'écouter}: consulter le schéma et le contenu de la base des immeubles.

\vfill
\begin{blocgauche}
\alert{Ces diapositives correspondent au support en ligne disponible sur
le site \url{http://sql.bdpedia.fr/}}
\end{blocgauche}

\end{frame}


\begin{frame}[fragile]{Valeurs nulles}

Une \alert{valeur nulle}, ou plus précisément \alert{valeur à \texttt{null}} est une
\alert{valeur manquante}. Ne pas confondre avec la valeur "null" ou "".

\vfill

Dans notre table des occupants, le prénom de Prof est à null.

\vfill
\begin{tabular}{c|c|c|c|c}
\textbf{id} & \textbf{prénom} & \textbf{nom} & \textbf{profession} & \textbf{idAppart}\\ \hline 
1 &  & Prof & Enseignant & 202\\ \hline 
2 & Alice & Grincheux & Cadre & 103\\ \hline 
3 & Léonie & Atchoum & Stagiaire & 100\\ \hline 
4 & Barnabé & Simplet & Acteur & 102\\ \hline 
5 & Alphonsine & Joyeux & Rentier & 201\\ \hline 
6 & Brandon & Timide & Rentier & 104\\ \hline 
7 & Don-Jean & Dormeur & Musicien & 200\\ \hline 
\end{tabular}

\vfill

\begin{blocgauche}
La présence de valeurs à \texttt{null} \alert{fausse} le résultat attendu des requêtes.
\end{blocgauche}
\end{frame}


\begin{frame}[fragile]{Comparaisons avec valeurs nulles}

On ne sait pas comparer un valeur manquante, ou lui appliquer une fonction.

\vfill


\begin{minted}[baselinestretch=.9,
               fontsize=\small,
               framesep=2mm]{sql}
    select annéeNaissance
     from Artiste as a
     where a.nom='Depardieu'

\end{minted}

Prof (pas de prénom) n'est pas trouvé.

\vfill

\begin{minted}[baselinestretch=.9,
               fontsize=\small,
               framesep=2mm]{sql}
select  * from Personne where prénom not like  '%'
\end{minted}
Prof n'est pas trouvé non plus!
\vfill

\begin{blocgauche}
Une comparaison avec \texttt{null} ne donne ni vrai, ni faux,
mais une troisième valeur de vérité, \texttt{unknown}. 
\end{blocgauche}
\end{frame}


%%%%%%%%
\begin{frame}[fragile]{Calculs avec valeurs à null}

Tout calcul appliqué avec une valeur à \texttt{null} renvoie \texttt{null} !
\vfill


\begin{minted}[baselinestretch=.9,
               fontsize=\small,
               framesep=2mm]{sql}
select concat(prénom, ' ', nom) as 'nomComplet'
from Personne
\end{minted}

\vfill

\begin{tabular}{c}
\textbf{nomComplet}\\ \hline 
null\\ \hline 
Alice Grincheux\\ \hline 
Léonie Atchoum\\ \hline 
Barnabé Simplet\\ \hline 
Alphonsine Joyeux\\ \hline 
Brandon Timide\\ \hline 
Don-Jean Dormeur\\ \hline 
\end{tabular}

\end{frame}

%%%%%%%%%%%%%%%%%%%%
\begin{frame}[fragile]{Le test \texttt{is null}}

Seule approche correcte: il faut tester explicitement l'absence de valeur avec \texttt{is null}.

\vfill
En SQL: 

\begin{minted}[baselinestretch=.9,
               fontsize=\small,
               framesep=2mm]{sql}
   select * from Personne
   where prénom like '%' 
   or prénom is null
\end{minted}

\vfill

\alert{Attention}\, le test \texttt{prénom = null} \alert{ne marche pas}.

\vfill

\begin{blocgauche}
\alert{Conclusion}: éviter autant que possible les valeurs à \texttt{null}
en les interdisant (dans le schéma).
\end{blocgauche}
\end{frame}

%%%%%%%%%%%%%%%%%%%%
\begin{frame}[fragile]{La jointure externe}

On veut la liste des appartements avec leurs occupants.
\vfill


\begin{minted}[baselinestretch=.9,
               fontsize=\small,
               framesep=2mm]{sql}
    select idImmeuble, no, niveau, surface, nom, prénom
    from  Appart as a, Personne as p
    where  p.idAppart=a.id
\end{minted}

\vfill

\begin{tabular}{c|c|c|c|c|c}
\textbf{idI} & \textbf{no} & \textbf{niveau} & \textbf{surface} & \textbf{nom} & \textbf{prénom}\\ \hline 
2 & 2 & 2 & 250 & Prof & null\\ \hline 
1 & 52 & 5 & 50 & Grincheux & Alice\\ \hline 
1 & 1 & 14 & 150 & Atchoum & Léonie\\ \hline 
1 & 51 & 2 & 200 & Simplet & Barnabé\\ \hline 
2 & 1 & 1 & 250 & Joyeux & Alphonsine\\ \hline 
1 & 43 & 3 & 75 & Timide & Brandon\\ \hline 
2 & 10 & 0 & 150 & Dormeur & Don-Jean\\ \hline 
\end{tabular}

\vfill

\begin{blocgauche}
\alert{Il manque l'appartement 34 qui n'a pas d'occupant}.
\end{blocgauche}
\end{frame}

%%%%%%%%%%%%%%%%%%%%
\begin{frame}[fragile]{L'opérateur \texttt{outer join}}

L'opérateur algébrique \texttt{outer join} 

\begin{redbullet}
   \item Renvoie tous les nuplets de la table directrice (celle de gauche)
   \item Associe à chaque nuplet un nuplet de la table de droite \alert{si un tel nuplet existe}
   \item Sinon, les attributs
provenant de la table de droite sont affichés à \texttt{null}
\end{redbullet}

\vfill


\begin{minted}[baselinestretch=.9,
               fontsize=\small,
               framesep=2mm]{sql}
select idImmeuble, no niveau, surface, nom, prénom
from  Appart as a left outer join Personne as p on (p.idAppart=a.id)
\end{minted}

\vfill
On obtient, en plus de la jointure standard: 

\vfill
\begin{tabular}{c|c|c|c|c}
\hline
1 & 34 & 50 & null & null\\ \hline 
\end{tabular}

\end{frame}



%%%%%%%%%%%%%%%%%%%%
\begin{frame}[fragile]{Le tri, \texttt{order by}}

On peut demander explicitement le tri du résultat sur
une ou plusieurs expressions avec la clause \texttt{order by}

\vfill


\begin{minted}[baselinestretch=.9,
               fontsize=\small,
               framesep=2mm]{sql}
   select *  
   from Appart 
   order by surface, niveau
\end{minted}

\vfill

En ajoutant des clauses sur l'ordre du tri (\texttt{ascending} ou \texttt{descending})

\vfill

\begin{minted}[baselinestretch=.9,
               fontsize=\small,
               framesep=2mm]{sql}
    select *  
    from Appart 
    order by surface desc, niveau desc
\end{minted}

\end{frame}
%%%%%%%%%%%
\begin{frame}{À retenir}

Les fondements du langage SQL (logique ou algèbre) sont une base sur laquelle
beaucoup d'extensions pratiques sont possibles.

\vfill

\begin{redbullet}
   \item Valeurs à \texttt{null}, jointures externes, tris
  \item  Mais aussi des fonctions, spécifiques à chaque système
  \item Ne change en rien \alert{l'interprétation} du langage, que vous devez maintenant maîtriser.
\end{redbullet}

\vfill
\begin{blocgauche}
SQL est un standard, mais chaque système propose des spécificités au-delà du standard.
\end{blocgauche}
\end{frame}

